\documentclass[a4paper]{article}

% Set margins
\usepackage[hmargin=2.5cm, vmargin=3cm]{geometry}

\frenchspacing

% Language packages
\usepackage[utf8]{inputenc}
\usepackage[T1]{fontenc}
\usepackage[magyar]{babel}

% AMS
\usepackage{amssymb,amsmath}

% Graphic packages
\usepackage{graphicx}

% Colors
\usepackage{color}
\usepackage[usenames,dvipsnames]{xcolor}

% Listings
\usepackage{listings}

% Question
\newenvironment{question}[1]
{\noindent\textcolor{OliveGreen}{$\circ$ \textit{#1}}

\smallskip

\color{Gray}

}{\bigskip}

% Task
\newenvironment{task}[1]
{\noindent\textcolor{RoyalBlue}{$\circ$ \textit{#1}}

\smallskip

\color{Gray}

}{\bigskip}

% Notification
\newenvironment{notification}[1]
{\noindent\textcolor{Peach}{$\circ$ \textit{#1}}

\smallskip

\color{Gray}

}{\bigskip}

% Problem
\newenvironment{problem}[1]
{\noindent\textcolor{OrangeRed}{$\circ$ \textit{#1}}

\smallskip

\color{Gray}

}{\bigskip}

% Solution
\newenvironment{solution}
{\color{RoyalBlue}}{\bigskip}

% Comment
\newenvironment{comment}
{\color{Green}}{\bigskip}

\lstset{% setup listings
        framexleftmargin=16pt,
        framextopmargin=6pt,
        framexbottommargin=6pt, 
        frame=single, rulecolor=\color{black}
}

\begin{document}

\begin{center}
   \large \textbf{Valószínűségszámítás - típusfeladatok - 2. ZH}
\end{center}

\vskip 1cm

\section{Nevezetes eloszlások}

\subsection{Poisson eloszlás}

\noindent \textbf{Típusok}

\begin{itemize}
\item Adott a $\lambda$ paraméter, és ki kell számítani valamilyen valószínűséget.
\item Adott egy valószínűség, és ki kell számítani a $\lambda$ paramétert.
\item Adott a $\lambda$ és egy valószínűség, és ki kell számolni egy hiányzó paramétert.
\item Darabszámok intervallumát vagy több véletlen kísérletet (mintavételt) kell vizsgálni (például a mazsolás feladatnál, hogy legalább 2 szeletben nincs mazsola).
\end{itemize}

\subsection{Exponenciális eloszlás}

\noindent \textbf{Típusok}

\begin{itemize}
\item Adott a $\lambda$ paraméter és egy valószínűséget kell meghatározni.
\item Adott valószínűség, és a $\lambda$ a meghatározandó.
\item A tartományt magát kell meghatározni (például, hogy csak az egyik vége adott, és mennyi lesz a másik).
\end{itemize}

\subsection{Nevezetes eloszlások kapcsolata}

\textbf{Típusok}

\begin{itemize}
\item Intervallumba esés valószínűsége.
\item Intervallum meghatározása valószínűség alapján.
\item Hiányzó várható érték vagy szórás paraméter.
\item Binomiális eloszlás közelítése normáis eloszlással (Moivre-Laplace tétel)
\end{itemize}

\section{Normális eloszlás}

\noindent \textbf{Típusok}

\begin{itemize}
\item Adott a várhatóérték és a szórásnégyzet, és ki kell számolni valamilyen valószínűséget.
\item Adott egy valószínűség, és ki kell számolni a várható értéket.
\item Adott valószínűség esetén a hozzá tartozó intervallumot kell meghatározni.
\end{itemize}

\noindent \textbf{Észrevételek}

\begin{itemize}
\item A várható érték az első momentum, a szórásnégyzet pedig a második. Előfordulhat, hogy ilyen formában lesz megfogalmazva a feladat.
\item A feladat megoldásához normális eloszlás táblázatot célszerű használni.
\end{itemize}

\section{Egyenlőtlenségek}

\subsection{Markov-egyenlőtlenség}

\noindent \textbf{Típusok}

\begin{itemize}
\item Adott a várható érték, és be kell látni, hogy bizonyos egyenlőtlenség teljesül.
\item A $c$ értéket kell meghatározni.
\item Az $E(\xi)$ értéket kell meghatározni.
\end{itemize}

\noindent \textbf{Észrevételek}

\begin{itemize}
\item A feladat szövegezéséből kell felismerni, hogy a Markov-egyenlőtlenségre lesz szükség. (Az egyenlőtlenség valamilyen formában szerepel benne, és a várható értéken kívül mást nem lehet tudni, vagy épp azt kell kiszámolni.)
\item Az egyenlőtlenséget a komplementer esemény alapján át kell tudni rendezni.
\end{itemize}

\subsection{Csebisev-egyenlőtlenség}

\noindent \textbf{Típusok}

\begin{itemize}
\item Ismert a szórásnégyzet. A várhatóértéktől való eltérés valószínűségére kell becslést adni.
\item Ismert a várhatóértéktől való eltérés valószínűsége. A szórásnégyzetre kell becslést adni.
\item Ismert az egyenlőtlenség jobb oldala és a szórásnégyzet. A várhatóértéktől való eltérésre kell becslést adni.
\end{itemize}

\noindent \textbf{Észrevételek}

\begin{itemize}
\item Ennél is szükség lehet az egyenlőtlenség átrendezésére.
\item A várható érték is meg lehet adva.
\item A szórásnégyzet előfordulhat, hogy ismert eloszlásból számítandó ki.
\end{itemize}

\subsection{Bernoulli-féle nagy számok törvénye}

\noindent \textbf{Típusok}

\begin{itemize}
\item Valamit gyártunk/készítünk, és szeretnénk minél kisebb intervallumot meghatározni a gyártás pontosságára vonatkozóan. (Például, hogy ha azt mondjuk, hogy a legyártott termék hossza $\pm 1$ cm legalább 0.95 valószínűséggel az jobb, mint ha azt mondjuk, hogy $\pm 200$ cm.)
\item Ismert a pontosság, és meg szeretnénk határozni, hogy az milyen valószínűséggel jelenthető ki.
\item A szükséges kísérletek számát szeretnénk becsülni.
\item A vizsgált esemény bekövetkezésének valószínűségét szeretnénk vizsgálni.
\end{itemize}

\noindent \textbf{Észrevételek}

\begin{itemize}
\item Az alapján lehet felismerni, hogy erre az egyenlőtlenségre van szükség, hogy Bernoulli-kísérletsorozatról van szó, a relatív gyakoriság valószínűségtől való eltérését kell becsülni, és hogy jellemzően a kísérletek száma is adott.
\end{itemize}

\section{Véletlen vektorok - Diszkrét eset}

\noindent \textbf{Típusok}

\begin{itemize}
\item Adott egy eloszlás táblázatos formában. Fel kell írni és rajzolni az eloszlásfüggvényét.
\item Ki kell számítani a peremeloszlásokat (esetleg azokat is ábrázolni kell).
\item Ki kell számítani adott intervallumba esés valószínűségét. (Meg lehet adva szövegesen, vagy közvetlenül $P$ függvénnyel felírva.)
\item Ki kell számítani a feltételes valószínűségeket.
\item Kiszámítandó az $E(\xi), D^2(\xi), D(\xi), E(\eta), D^2(\eta), D(\eta)$ értéke.
\item Ki kell számítani a $\xi$-nek az $\eta$-ra vonatkozó, illetve az $\eta$-nak a $\xi$-re vonatkozó regressziós függvényét.
\end{itemize}

\section{Véletlen vektorok - Folytonos eset}

\noindent \textbf{Típusok}

\begin{itemize}
\item Adott az együttes sűrűségfüggvény és fel kell írni az eloszlásfüggvényt.
\item Adott az együttes eloszlásfüggvény, és meg kell határozni az együttes sűrűségfüggvényt.
\item Fel kell rajzolni tartomány formájában ("felülről nézve") a sűrűségfüggvényt.
\item Ki kell számítani az $E(\xi), D^2(\xi), D(\xi), E(\eta), D^2(\eta), D(\eta)$ értékeket.
\item Meg kell határozni a $\xi$ vagy az $\eta$ mediánját ($\text{med}\{\xi\}, \text{med}\{\eta\}$).
\item Fel kell írni a feltételes sűrűségfüggvényeket.
\item Ki kell számítani az intervallumba esés valószínűségét.
\item Meg kell határozni a sűrűségfüggvény egy hiányzó paraméterét.
\item Meg kell vizsgálni, hogy $\xi$ és $\eta$ függetlenek-e.
\end{itemize}

\noindent \textbf{Észrevételek}

\begin{itemize}
\item A hiányzó paraméter lehet a tartomány megadásában is.
\item Akkor is érdemes ábrázolni, hogy ha az feladatként nem szerepel, mert gyakran abból is leolvashatóak az eredmények, és ellenőrízhetők.
\end{itemize}

\section{Leíró statisztikák}

\noindent \textbf{Típusok}

\begin{itemize}
\item Adott minta alapján a következőket kell kiszámítani: átlag, medián, MAD, korrigált tapasztalati szórás, kvartilisek, $p$-kvantilisek számítása
\item Fel kell rajzolni a tapasztalati eloszlásfüggvényt, a gyakoriság vagy a sűrűséghisztogramot.
\end{itemize}

\end{document}
